\subsection{Deleting Names}

Assignment creates or changes a binding between a name and an object. The statement \texttt{del} can remove such a binding.

The simplest form is:

\begin{verbatim}
del name
\end{verbatim}

This does not mean: destroy the object immediately. It means: remove this name from the current namespace.

Start with an ordinary binding:

\begin{verbatim}
quote = "No soup for you."

print(f"quote: {quote}")
print(f"type(quote): {type(quote)}")
print(f"id(quote):   {id(quote)}")
\end{verbatim}

Possible output:

\begin{verbatim}
quote: No soup for you.
type(quote): <class 'str'>
id(quote):   2168389812912
\end{verbatim}

The name \texttt{quote} refers to a string object. After \texttt{del quote}, that name no longer exists in the current namespace.

\subsubsection{Removing a Binding with \texttt{del}}

The statement \texttt{del} removes a name binding.

\begin{verbatim}
line = "Nobody expects the Spanish Inquisition!"

print("before del")
print(f"line: {line}")
print(f"type(line): {type(line)}")

del line

print("after del")
print(line)
\end{verbatim}

This program does not finish normally. After \texttt{del line}, the name \texttt{line} is no longer bound. Trying to read it raises a \texttt{NameError}.

A safer version catches the error:

\begin{verbatim}
line = "Nobody expects the Spanish Inquisition!"

print("before del")
print(f"line: {line}")
print(f"type(line): {type(line)}")

del line

print()
print("after del")

try:
    print(line)
except NameError as err:
    print("name lookup failed")
    print(f"error object: {err}")
    print(f"type(err):    {type(err)}")

    print()
    print("selected NameError methods/attributes:")
    print(f"{'args':<12} {'args' in dir(err)}")
    print(f"{'add_note':<12} {'add_note' in dir(err)}")
    print(f"{'with_traceback':<12} {'with_traceback' in dir(err)}")
\end{verbatim}

Possible output:

\begin{verbatim}
before del
line: Nobody expects the Spanish Inquisition!
type(line): <class 'str'>

after del
name lookup failed
error object: name 'line' is not defined
type(err):    <class 'NameError'>

selected NameError methods/attributes:
args         True
add_note     True
with_traceback True
\end{verbatim}

The important point is not the exception syntax. The important point is the binding change:

\begin{verbatim}
before:
    line ---> str object "Nobody expects the Spanish Inquisition!"

after:
    line no longer exists as a usable name here
\end{verbatim}

The object and the name are different things. \texttt{del line} deletes the binding named \texttt{line}. It does not mean that the programmer has directly reached into memory and destroyed a string object.

A name can also be rebound after deletion:

\begin{verbatim}
fact = "Chuck Norris counted to infinity. Twice."

print(f"fact: {fact}")
print(f"type(fact): {type(fact)}")

del fact

fact = 42

print()
print("after deletion and new assignment")
print(f"fact: {fact}")
print(f"type(fact): {type(fact)}")
\end{verbatim}

Possible output:

\begin{verbatim}
fact: Chuck Norris counted to infinity. Twice.
type(fact): <class 'str'>

after deletion and new assignment
fact: 42
type(fact): <class 'int'>
\end{verbatim}

This is a new binding using the same source-code name. The old binding was removed. Then assignment created a new binding.

Conceptually:

\begin{verbatim}
step 1:
    fact ---> str object "Chuck Norris counted to infinity. Twice."

step 2:
    del fact
    fact is not bound

step 3:
    fact ---> int object 42
\end{verbatim}

The word \texttt{del} is therefore about removing a reference from a namespace. It is not the same operation as setting a name to \texttt{None}.

Compare:

\begin{verbatim}
value = "D'oh!"
value = None

print(f"value: {value}")
print(f"type(value): {type(value)}")
\end{verbatim}

Possible output:

\begin{verbatim}
value: None
type(value): <class 'NoneType'>
\end{verbatim}

Here the name still exists. It is bound to the singleton object \texttt{None}.

With \texttt{del}, the name itself is removed:

\begin{verbatim}
value = "D'oh!"
del value

try:
    print(value)
except NameError as err:
    print(f"cannot read value: {err}")
\end{verbatim}

Possible output:

\begin{verbatim}
cannot read value: name 'value' is not defined
\end{verbatim}

So the difference is:

\begin{verbatim}
value = None   -> name exists and refers to None
del value      -> name no longer exists here
\end{verbatim}

For ordinary beginner code, \texttt{del} is not used often with simple variables. It is more common later with container elements, module cleanup, or cases where the programmer deliberately wants to remove a name.

\subsubsection{Name Deletion, Object Reachability, and Possible Reference Count Changes}

Deleting one name does not necessarily make the object unreachable. If another name still refers to the same object, the object is still usable.

\begin{verbatim}
a = "It is merely a flesh wound."
b = a

print("before del a")
print(f"a: {a}")
print(f"b: {b}")
print(f"id(a): {id(a)}")
print(f"id(b): {id(b)}")
print(f"a is b: {a is b}")

del a

print()
print("after del a")
print(f"b: {b}")
print(f"id(b): {id(b)}")
\end{verbatim}

Possible output:

\begin{verbatim}
before del a
a: It is merely a flesh wound.
b: It is merely a flesh wound.
id(a): 2168389813296
id(b): 2168389813296
a is b: True

after del a
b: It is merely a flesh wound.
id(b): 2168389813296
\end{verbatim}

Before deletion:

\begin{verbatim}
a ----\
      ---> str object "It is merely a flesh wound."
b ----/
\end{verbatim}

After \texttt{del a}:

\begin{verbatim}
b ---> str object "It is merely a flesh wound."
\end{verbatim}

The object is still reachable through \texttt{b}. Therefore, it cannot be treated as gone.

This is the key rule:

\begin{quote}
An object can remain alive as long as it is reachable from at least one active reference.
\end{quote}

In CPython, the main Python implementation, object lifetime is strongly connected to reference counting. A reference count is the number of active references that CPython is currently tracking for an object. When a binding is removed, the object's reference count may decrease.

This can be observed with the standard \texttt{sys} module:

\begin{verbatim}
import sys

quote = "These pretzels are making me thirsty."

print(f"quote: {quote}")
print(f"type(quote): {type(quote)}")

print()
print(f"refcount through quote: {sys.getrefcount(quote)}")

alias = quote

print()
print("after alias = quote")
print(f"quote is alias: {quote is alias}")
print(f"refcount through quote: {sys.getrefcount(quote)}")

del alias

print()
print("after del alias")
print(f"refcount through quote: {sys.getrefcount(quote)}")
\end{verbatim}

Possible output:

\begin{verbatim}
quote: These pretzels are making me thirsty.
type(quote): <class 'str'>

refcount through quote: 4

after alias = quote
quote is alias: True
refcount through quote: 5

after del alias
refcount through quote: 4
\end{verbatim}

The exact numbers should not be memorized. They depend on interpreter details, temporary references, constants, and how the code is executed. Also, \texttt{sys.getrefcount()} itself temporarily receives the object as an argument, so the displayed count includes that temporary reference.

The useful observation is the direction:

\begin{itemize}
    \item creating another binding can increase the reference count;
    \item deleting that binding can decrease the reference count;
    \item the object remains reachable while at least one usable reference remains.
\end{itemize}

A clearer example uses a list because lists are normally not reused as constants in the same way small integers or some strings may be.

\begin{verbatim}
import sys

items = ["spam", "eggs", 42]

print(f"items: {items}")
print(f"type(items): {type(items)}")

print("selected list methods:")
print(f"{'append':<12} {'append' in dir(items)}")
print(f"{'extend':<12} {'extend' in dir(items)}")
print(f"{'pop':<12} {'pop' in dir(items)}")

print()
print(f"initial refcount: {sys.getrefcount(items)}")

same_items = items

print()
print("after same_items = items")
print(f"items is same_items: {items is same_items}")
print(f"refcount: {sys.getrefcount(items)}")

del same_items

print()
print("after del same_items")
print(f"refcount: {sys.getrefcount(items)}")
\end{verbatim}

Possible output:

\begin{verbatim}
items: ['spam', 'eggs', 42]
type(items): <class 'list'>
selected list methods:
append       True
extend       True
pop          True

initial refcount: 2

after same_items = items
items is same_items: True
refcount: 3

after del same_items
refcount: 2
\end{verbatim}

Again, the exact count is not the lesson. The lesson is that bindings participate in reachability.

If the last reachable name is deleted, the program can no longer access the object through that name:

\begin{verbatim}
items = ["spam", "eggs", 42]

print(f"items before del: {items}")

del items

try:
    print(items)
except NameError as err:
    print("items is no longer a usable name")
    print(f"message: {err}")
\end{verbatim}

Possible output:

\begin{verbatim}
items before del: ['spam', 'eggs', 42]
items is no longer a usable name
message: name 'items' is not defined
\end{verbatim}

At that point, if no other references exist, CPython can usually reclaim the object immediately by reference counting. Other Python implementations may use different memory-management strategies. Even in CPython, object lifetime can be affected by reference cycles, temporary references, exception traces, interned constants, and implementation details.

Therefore the careful rule is:

\begin{quote}
\texttt{del} removes a binding. Removing a binding may reduce an object's reference count. If the object becomes unreachable, the interpreter may reclaim it. The programmer should not read \texttt{del name} as a direct memory-free instruction.
\end{quote}

For this chapter, the practical mental model is enough:

\begin{verbatim}
assignment  -> create or change a name-to-object binding
reassignment -> move a name to another object
a = b       -> bind another name to the same object
del name    -> remove one name-to-object binding
\end{verbatim}